\documentclass[12pt, b5paper]{article}
\usepackage{ctex}

\usepackage{geometry}
\geometry{b5paper,left=2cm,right=1cm,top=2.5cm,bottom=2.5cm}
\usepackage{chemformula} %化学公式宏包
\usepackage[version=4]{mhchem} %化学公式宏包
\usepackage{siunitx} %数字，单位宏包
\sisetup{
	separate-uncertainty = true,
	inter-unit-product = \ensuremath{{}\cdot{}}
} %单位间加小圆点
\usepackage{tikz} %Tikz绘图包
\usetikzlibrary{cd}
\usetikzlibrary{chains}
\usetikzlibrary{arrows}
\usetikzlibrary{arrows.meta} %画箭头
\usetikzlibrary{commutative-diagrams} %交换图
\usetikzlibrary{positioning,automata}
\usepackage[all]{xy}
\usepackage{booktabs} %三线表
\usepackage{multirow} %合并单元格
\usepackage{makecell} %表格内强制换行
\usepackage{array}
\usepackage{booktabs} %控制表格行高
\usepackage{colortbl}  %彩色表格需要加载的宏包
\usepackage{amsmath}
\usepackage{unicode-math}
\usepackage{fontspec} %字体
\newcommand*{\cred}{\textcolor{red}}
\setmainfont{Times New Roman}
\setmathfont{XITS Math}
\setCJKmainfont{SimSun}[AutoFakeBold, AutoFakeSlant] 

\setlength{\parindent}{0pt} %无缩进
\usepackage{setspace}
\usepackage{fancyhdr} %设置页码
\usepackage{lastpage} %获取总页数
\pagestyle{fancy} %fancyhdr宏包新增的页面风格
\renewcommand\headrulewidth{0pt} %取消页眉中的装饰分割线
\fancyhf{}
\cfoot{\footnotesize 第 \thepage\ 页(共 \pageref{LastPage}页)}%当前页 of 总页数



\begin{document}


\begin{spacing}{1.625}

\centerline{{\huge \bf{河北农业大学课程考试试卷}}}

\centerline{\bf{\large 20\underline{25}年春（\quad ）/秋（\ \surd \ ）学期}}

{\bf{考试科目：\underline{\quad 化学反应工程 \quad} \hfill 姓名：\underline{ \qquad \qquad} \hfill 学号：\underline{ \qquad \qquad \qquad \qquad}}}

{\bf{学院：\underline{\quad 理工系\quad } \hfill 专业班级：\underline{\qquad \qquad \qquad \qquad} \hfill 卷别：\cred{B} \hfill 考核方式：\cred{开卷}}}

（注：考生务必将答案写在答题纸上，标清楚大小题号，写在本试卷上无效）

\bigskip


{\bf{一、简答题，用一句话回答（每题5分，共30分）}}

1. （5分） 工业上选择反应器类型时，需要考虑哪些方面？（说出一个即可）

\bigskip

2. （5分）在反应器设计中，“体积”是一个关键参数。请说出一个影响反应器体积大小的化学因素。

\bigskip

3. （5分）实验室的小反应器放大到工业大装置时，可能会遇到什么问题？

\bigskip

4. （5分）请说出一个全混流反应器的显著特点。

\bigskip

5. （5分）请写出两种常见的理想反应器模型名称。

\bigskip

6. （5分）如果把一个大颗粒的催化剂磨成小颗粒，主要会改善内扩散过程还是外扩散过程？

\bigskip

{\bf{二、问答题（每题10分，共30分）}}

1. （10分）在多相催化反应中，“内扩散”和“外扩散”的区别是什么？给出消除内、外扩散阻力的方法。

\bigskip

2. （10分）为什么说化学反应工程是化学工程的“核心”课程？

\bigskip

3. （10分）请画出全混流反应器和平推流反应器的反应物浓度随反应器位置变化的示意图。

\begin{tikzpicture}[>=Stealth] % 使用更漂亮的箭头样式

  % 绘制坐标轴，只保留第一象限
  \draw[->, line width=1pt] (0,0) -- (4,0) node[below] {出口}; % x轴
  \draw[->, line width=1pt] (0,0) -- (0,4) node[left] {$c$}; % y轴

  \node[below left] at (0.5,0) {入口};

  \draw[->, line width=1pt] (8,0) -- (12,0) node[below] {出口}; % x轴
  \draw[->, line width=1pt] (8,0) -- (8,4) node[left] {$c$}; % y轴

  \node[below left] at (8.5,0) {入口};

\end{tikzpicture}

\bigskip

{\bf{二、计算题(共40分)}}

1. (40分)

在等温下进行液相反应\ch{A -> B}，在该条件下的反应速率方程为
\[
    r_\mathrm{A} = 0.8c_\mathrm{A}^2 \qquad [\si{mol.L^{-1}.h^{-1}}]
\]
原料A的初始浓度为3 mol/L。

(1) 求反应4 h时A的转化率。

(2) 每天处理\SI{24}{m^3}原料，要求最终转化率为80\%，反应在连续釜式反应器中进行，求反应体积。

(3) 条件同上，反应在活塞流管式反应器中进行，求反应体积。

(4) 条件同上，反应在间歇釜式反应器中进行，辅助时间为0.5 h，求反应体积。


\end{spacing}
\end{document}
